\mxunit{example}{Exemple d'unité}
\begin{objectives}
Comprendre la structure d'une unité ; utiliser les environnements canoniques.
\end{objectives}
\section{Une définition}
\begin{definition}[Suite arithmétique]
Une suite $(u_n)$ est arithmétique de raison $r$ si, pour tout $n\in\N$, $u_{n+1}=u_n+r$.
\end{definition}
\begin{example}
Avec $u_0=2$ et $r=0{,}5$, on obtient $u_3=3{,}5$. \mxchip{5 min}
\end{example}
\begin{pitfall}
La raison se lit sur la différence $u_{n+1}-u_n$, pas sur le quotient.
\end{pitfall}
\begin{theorem}
Pour tout $n\in\N$, $u_n=u_0+nr$.
\end{theorem}
\begin{proof}
Par récurrence sur $n$ : l'égalité est vraie pour $n=0$, et si $u_n=u_0+nr$ alors $u_{n+1}=u_0+(n+1)r$.
\end{proof}
\begin{exercise}[Calcul de termes]
Calculer $u_{10}$ pour $u_0=-3$ et $r=2$. Réponse : $u_{10}=\mxblank{17}$.
\mxlines{2}
\end{exercise}
\begin{solution}
$u_{10}=-3+10\times2=17$.
\end{solution}
\begin{figurebox}[Les termes sont alignés.]
\begin{tikzpicture}\begin{axis}[mxaxis,width=7cm,height=4cm,xmin=0,xmax=5,ymin=0,ymax=5]
\addplot[mxcurve,only marks] coordinates {(0,2)(1,2.5)(2,3)(3,3.5)(4,4)};
\end{axis}\end{tikzpicture}
\end{figurebox}
\begin{teachernote}
Le \mxdate{2026-09-01}, faire reformuler la définition à l'oral.
\end{teachernote}
